% Source: notes/2026-10-21-agnostic-learning.jpg

\begin{frame}{Agnostic learning}
  \textcolor{muted}{As distinct from the realizable setting.}

  \vspace{1em}
  Data
  \[
    \big((x_1,y_1),\dots,(x_n,y_n)\big) \sim \D^n
  \]

  Want a classifier $h$ with low error
  \[
    \err(h) = \Prob\big[h(x) \neq y\big].
  \]

  \pause
  \vspace{0.5em}
  \alert{This may be impossible, or very hard.}
\end{frame}

\begin{frame}{Agnostic learning: a relative guarantee}
  Data $\big((x_1,y_1),\dots,(x_n,y_n)\big) \sim \D^n$

  Comparator class $\F$ of functions $f : \X \to \{0,1\}$

  \vspace{1em}
  With probability at least $1-\delta$,
  \[
    \err(h) \;\le\; \min_{f \in \F} \err(f) \;+\; \varepsilon .
  \]

  \vspace{0.5em}
  {\color{muted}The excess-risk framework of statistical learning theory.}
\end{frame}

\begin{frame}{Agnostic learning: predicting the truth}
  Learn a predictor $p : \X \to [0,1]$.

  \vspace{0.3em}
  The holy grail:
  \[
    p(x) = \E[\,y \mid x\,]
  \]
  {\color{muted}Learn the observed data up to irreducible noise.}

  \pause
  \vspace{0.8em}
  We will learn some $\hat p$:
  \[
    \begin{array}{rcl}
      d(p,\hat p) & \longrightarrow & \text{epistemic uncertainty}\\[0.4em]
      d(\hat p, y) & \longrightarrow & \text{aleatoric uncertainty}
    \end{array}
  \]
\end{frame}
